% Generated draft: requires mathematical and source review.
\documentclass[11pt]{article}
\usepackage[margin=25mm]{geometry}
\usepackage{fontspec}
\setmainfont{Latin Modern Roman}
\usepackage{amsmath,amssymb,mathtools,unicode-math}
\setmathfont{Latin Modern Math}
\usepackage{enumitem,needspace,xurl,hyperref}
\setlength{\parindent}{0pt}
\setlength{\parskip}{5pt}
\setlength{\emergencystretch}{3em}
\newcommand{\N}{\mathbb{N}}
\newcommand{\Z}{\mathbb{Z}}
\newcommand{\Q}{\mathbb{Q}}
\newcommand{\R}{\mathbb{R}}
\newcommand{\C}{\mathbb{C}}
\newcommand{\F}{\mathbb{F}}
\newcommand{\A}{\mathbb{A}}
\newcommand{\PP}{\mathbb{P}}
\newcommand{\E}{\mathbb{E}}
\newcommand{\eps}{\varepsilon}
\newcommand{\dd}{\,\mathrm{d}}
\begin{document}
\section*{UnsolvedMath ID: 701538}
\textbf{Source problem code:} catalogue entry
\subsection{Definitions and mathematical statement}
Let \(p\) be a prime. A finite group is \(p\)-soluble if it has a normal series whose factors are either \(p\)-groups or groups of order coprime to \(p\). A group variety \(\mathfrak V\) is a class of groups defined by group identities; it is proper if it is not the variety of all groups.

For a finite group \(G\), consider normal series
\[
1=G_0\trianglelefteq G_1\trianglelefteq\cdots\trianglelefteq G_m=G,
\qquad G_i\trianglelefteq G,
\]
whose nontrivial factors are either \(p\)-soluble or finite direct products of nonabelian simple groups of order divisible by \(p\). The supplied source defines the \emph{non-\(p\)-soluble length} \(\lambda_p(G)\) as the number of non-\(p\)-soluble factors in a shortest such series. We retain that wording: the supplied record does not explicitly distinguish minimizing the total number of factors from minimizing the number of non-\(p\)-soluble factors, and this convention requires confirmation against the original source.

\textbf{Problem (Khukhro--Shumyatsky, Kourovka Notebook 18.110).} Is it true that for every prime \(p\) and every proper group variety \(\mathfrak V\), there exists a nonnegative integer \(B=B(p,\mathfrak V)\) such that every finite group \(G\) satisfies
\[
\bigl(P\in\mathfrak V\text{ for a Sylow }p\text{-subgroup }P\text{ of }G\bigr)
\quad\Longrightarrow\quad
\lambda_p(G)\le B?
\]
Membership in \(\mathfrak V\) for one Sylow \(p\)-subgroup is equivalent to membership for every Sylow \(p\)-subgroup, since they are conjugate. The bound may depend on \(p\) and \(\mathfrak V\), but not on \(G\).

\textbf{Supplied qualification.} The record labels the problem open, but its \emph{Comment of 2021} reports that the existence of such a bound for \(p=2\) was proved by F. Fumagalli, F. Leinen, and O. Puglisi, citing a publication dated 2022. This reported special case and the overall status label are supplied source information, not independently verified literature assessments.
\subsection{Short English statement}
For each fixed prime \(p\) and proper variety \(\mathfrak V\), must finite groups whose Sylow \(p\)-subgroups lie in \(\mathfrak V\) have uniformly bounded non-\(p\)-soluble length? The source reports an affirmative result for \(p=2\). Its wording about a shortest normal series is retained with the convention issue stated explicitly.
\subsection{Sources}
\begin{itemize}
\item Local UnsolvedMath registry: permanent ID \(701538\), exact source code \texttt{catalogue entry}, title \emph{Kourovka Notebook Problem 18.110}. Supplied upstream revision: \texttt{d371ee917ac9647fbb06cd6fc529deec6f3dda82}. The record was created and updated on 2026-10-16 and carries the supplied status \texttt{open}.
\item \emph{The Kourovka Notebook (Unsolved Problems in Group Theory)}, problem 18.110, proposed by E. I. Khukhro and P. Shumyatsky in 2014. The supplied background identifies the 18th issue (2014) and the 21st issue edition (2026), \texttt{arXiv:1401.0300v48}.
\item The supplied \emph{Comment of 2021} attributes the affirmative \(p=2\) case to F. Fumagalli, F. Leinen, and O. Puglisi, \emph{Proc. London Math. Soc. (3)}, \textbf{125}, no. 5 (2022), 1066--1082. No article title or link is supplied. This reference and its reported conclusion have not been independently checked.
\end{itemize}
\end{document}
